\section{Evaluation}
\label{sec:evaluation}


This section evaluates CRME through a scenario-based validation
framework. The evaluation focuses on four research questions related
to memory continuity, semantic organization, provenance preservation,
and research workflow support.



\subsection{Research Questions}


The evaluation is guided by the following research questions:


\textbf{RQ1: Memory Continuity.}
Can CRME preserve and recover research knowledge across separated
research sessions?


\textbf{RQ2: Semantic Organization.}
Can CRME represent relationships between research elements beyond
traditional document-based storage?


\textbf{RQ3: Provenance Preservation.}
Can CRME maintain traceable evolution paths for research decisions,
artifacts, and results?


\textbf{RQ4: Research Workflow Support.}
Can CRME support continuous research activities from initial ideas
toward scientific outputs?



% =====================================================
\subsection{Experimental Environment}
% =====================================================


The evaluation was conducted using a controlled research workflow
scenario containing multiple research sessions, project updates,
decisions, and generated artifacts.

The objective of this evaluation is to validate the architectural
capabilities of CRME rather than measure general artificial intelligence
performance.


\begin{table}[t]
\caption{Experimental Environment}
\label{tab:experimental_environment}

\centering
\small

\begin{tabular}{ll}

\toprule

Component & Configuration \\

\midrule

Operating System & Linux / Termux Environment \\

Implementation & Python 3.x \\

Storage Model & JSON-based Persistent Memory \\

Architecture & Modular Research Framework \\

Evaluation Type & Scenario-based Validation \\

\bottomrule

\end{tabular}

\end{table}



% =====================================================
\subsection{Evaluation Metrics}
% =====================================================


The following evaluation dimensions were considered:


\begin{itemize}

\item
\textbf{Memory Persistence:}
Ability to preserve research information beyond individual sessions.


\item
\textbf{Semantic Connectivity:}
Ability to represent relationships between research objects.


\item
\textbf{Provenance Completeness:}
Ability to track the origin and evolution of research information.


\item
\textbf{Workflow Continuity:}
Ability to support future research activities using previous knowledge.


\end{itemize}




% =====================================================
\subsection{Baseline Comparison}
% =====================================================


Table~\ref{tab:baseline_comparison} compares CRME with representative
research knowledge management approaches.


\begin{table}[t]

\caption{Capability Comparison Between CRME and Existing Approaches}

\label{tab:baseline_comparison}


\centering
\small


\begin{tabular}{lcccc}

\toprule

Approach &
Memory &
Semantic &
Provenance &
Continuity
\\

\midrule


Document Systems &
Limited &
No &
No &
No
\\


Workflow Systems &
Partial &
Partial &
Yes &
Limited
\\


Knowledge Graph Systems &
Partial &
Yes &
Partial &
Limited
\\


Retrieval-Augmented Systems &
Partial &
Yes &
Limited &
Limited
\\


CRME &
Yes &
Yes &
Yes &
Yes
\\


\bottomrule

\end{tabular}


\end{table}




% =====================================================
\subsection{Ablation Study}
% =====================================================


To analyze the contribution of individual CRME components,
an architectural ablation analysis was performed by removing
major modules.


The evaluated configurations include:


\begin{itemize}


\item
\textbf{Full CRME:}
Complete architecture including memory, session, project,
and semantic components.


\item
\textbf{Without Semantic Engine:}
Removes semantic relationship discovery capability.


\item
\textbf{Without Session Engine:}
Removes cross-session continuity management.


\item
\textbf{Without Project Engine:}
Removes project evolution tracking capability.


\end{itemize}



\begin{table}[t]

\caption{Ablation Analysis of CRME Components}

\label{tab:ablation_study}


\centering
\small


\begin{tabular}{lcccc}

\toprule

Configuration &
Memory &
Semantic &
Session &
Project
\\

\midrule


Full CRME &
Yes &
Yes &
Yes &
Yes
\\


Without Semantic Engine &
Yes &
No &
Yes &
Yes
\\


Without Session Engine &
Yes &
Yes &
No &
Yes
\\


Without Project Engine &
Yes &
Yes &
Yes &
No
\\


\bottomrule


\end{tabular}


\end{table}





% =====================================================
\subsection{Requirement Validation}
% =====================================================


Table~\ref{tab:validation_matrix} summarizes the relationship between
CRME requirements and implemented modules.


\begin{table}[t]

\caption{CRME Requirement Validation Matrix}

\label{tab:validation_matrix}


\centering
\small


\begin{tabular}{lll}

\toprule

Requirement &
Module &
Status
\\

\midrule


Persistent Memory &
Memory Engine &
Validated
\\


Session Recovery &
Session Engine &
Validated
\\


Project Evolution &
Project Engine &
Validated
\\


Semantic Linking &
Semantic Engine &
Validated
\\


Knowledge Querying &
Query Engine &
Validated
\\


\bottomrule


\end{tabular}


\end{table}




% =====================================================
\subsection{Scalability Analysis}
% =====================================================


The modular architecture of CRME enables incremental expansion
of research memory without modifying the underlying object model.

The current evaluation focuses on architectural scalability,
including modular separation, component independence, and future
support for quantitative benchmarking.



\begin{table}[t]

\caption{Architectural Scalability Analysis}

\label{tab:scalability_analysis}


\centering
\small


\begin{tabular}{ll}

\toprule

Dimension &
Observation
\\

\midrule


Memory Growth &
New research objects can be incrementally added
\\


Component Extension &
Independent engines enable modular expansion
\\


Knowledge Representation &
Object model supports additional research entities
\\


Future Benchmarking &
Retrieval accuracy and efficiency evaluation
\\


\bottomrule


\end{tabular}


\end{table}





% =====================================================
\subsection{Discussion}
% =====================================================


The evaluation demonstrates that CRME integrates capabilities that
are typically distributed across independent research management
systems.

The main contributions include:


\begin{itemize}


\item persistent research memory,

\item semantic organization,

\item provenance tracking,

\item session continuity,

\item project evolution management.


\end{itemize}


These properties enable a research environment where previous
knowledge can continuously influence future research activities.



% =====================================================
\subsection{Threats to Validity}
% =====================================================


Several limitations should be considered.


First, the current evaluation is based on controlled scenarios and
does not represent a large-scale user study.


Second, quantitative benchmarking using large research datasets is
required to measure retrieval accuracy and computational efficiency.


Third, semantic quality depends on the accuracy and completeness of
stored research representations.


These limitations define future experimental directions.



% =====================================================
\subsection{Reproducibility Statement}
% =====================================================


CRME follows a modular implementation approach.

The core components, including memory management, session management,
project management, semantic processing, and query mechanisms, are
implemented as independent modules.

The modular architecture enables future researchers to reproduce,
extend, and evaluate individual components separately.





